\chapter{Introduction}
\section{Introduction to Industrial Gearboxes}
	Industrial gear reducers are specialized mechanical assemblies designed to alter the rotational speeds and torque profiles 
    between a power source and driven industrial machinery. They serve two fundamental purposes: reducing the shaft velocity 
    and increasing the output torque proportionally. In practice, these sub-assemblies function as critical interface 
    nodes that connect different operating ends:
		\begin{itemize}
			\item \textbf{Prime Movers:} High-speed electric motors, internal combustion engines, or steam and fluid turbines.
			\item \textbf{Driven Machinery:} Heavy industrial plant equipment including material conveyors, pulverizing mills, 
			paper processing lines, commercial elevators, liquid mixers, and high-displacement screw pumps.
		\end{itemize}
	Standard industrial gearboxes are engineered for mainstream commercial infrastructure that demands reliable operational lifespans and conservative 
	safety margins. To distinguish them from high-performance drivetrains (such as aerospace propulsion or high-speed marine setups), their pitch line 
	velocity is usually restricted to under 25 m/s.

\section{Practical Challenges in Gearbox Engineering}
	To an engineer relying solely on textbook mechanics formulas, drafting a gearbox might seem trivial. However, developing a 
    reliable, manufacturable, and cost-effective gearbox requires managing highly coupled, non-linear failure modes. 
    Without practical experience, a designer can easily overlook systemic interactions, causing premature structural failures. 

    Developing a modern transmission requires balancing multiple coupled and often conflicting performance criteria simultaneously 
    for example: 
		\begin{itemize}
			\item \textbf{Weight vs. Efficiency and Packaging:}
			Thinning housing walls minimizes raw material costs and overall mass. 
			However, the resulting reduction in structural stiffness amplifies shaft misalignment. This misalignment triggers 
			uneven gear wear, edge loading, and higher power losses.
		\end{itemize}

\section{Traditional Gearbox Design}
	Drivetrain engineering relied on an empirical, handbook-driven methodology. This legacy pipeline leverages standardized physical manuals, historical layout 
	templates, and discrete look-up matrices to establish a mathematically viable gearbox layout. The workflow operates as a highly sensitive, sequential process 
	executed through manual calculation stages. 

	\subsection{System Input Metrics}
		Before executing any analytical calculations, the master input data matrix must be established from the application requirements. 

		A synthesis of cross-disciplinary aerospace programs, industrial engineering textbooks, and international gear standards (including NASA's A.R.T. methodologies, 
		AGMA 6013-B14, and ISO 6336) demonstrates a clear consensus on the foundational requirements for manual drivetrain synthesis. 

		Regardless of whether the application demands an ultra-lightweight helicopter 
		transmission or a heavy-duty industrial reduction drive, the absolute core 
		input variables that all engineering papers and design manuals agree upon are: 

		\subsubsection{Input Kinetic and Operational Boundaries}
			\begin{itemize}
				\item \textbf{P:} Nominal input power transmitted by the prime mover (expressed in kW or hp)
				\item \textbf{$n_{in}$:} Rotational speed of the high-speed input shaft in RPM
				\item \textbf{$n_{out}$:} Target rotational speed of the low-speed output shaft in RPM
				\item \textbf{$k_{a}$:} Empirical application service factor (accounts for prime mover and driven machine shock profiles)
				\item \textbf{$k_{t}$:} Lifecycle duty factor (normalizes operating hours per day against target lifespan)
			\end{itemize}

		\subsubsection{Material Mechanics Inputs}
			\begin{itemize}
				\item \textbf{$\rho$:} Mass density of the component material
				\item \textbf{E:} Modulus of Elasticity / Young's Modulus
				\item \textbf{$\nu$:} Poisson's ratio
				\item \textbf{$\sigma_b$:} Maximum allowable AGMA root bending fatigue stress
				\item \textbf{$\sigma_c$:} Maximum allowable AGMA surface compressive contact stress
				\item \textbf{$S_{y}$:} Tensile yield strength limit of the shaft material
				\item \textbf{$S_{n}$:} Fully corrected fatigue endurance limit of the shaft material
				\item \textbf{Lubrication and Thermal Limits:} Specific oil viscosity indexes, fluid properties, 
				and targeted operating temperatures required to prevent overheating and maintain oil-film thickness 
			\end{itemize}

		\subsubsection{Basic Geometric Setup Choices}
			\begin{itemize}
				\item \textbf{$\alpha_n$ Standardized Normal Pressure Angle:} The cutting profile orientation, 
				standardized internationally at ($20^\circ\text{C}$ or $25^\circ\text{C}$) that governs 
				contact force direction and controls tooth undercutting limits.
				\item \textbf{$\beta$ Helix Angle:} The angle of the gear teeth ($0^\circ\text{C}$) for spur gears; 
				typically $15^\circ\text{C}$ to $30^\circ\text{C}$ for helical setups, which introduces 
				axial thrust components into the load paths.
				\item \textbf{$\psi$ Face Width Factor:} The empirical ratio linking tooth width to module size, 
				dictating the width available to spread out mechanical loads. It is a designer-selected constraint 
				variable found by applying strict empirical lookup tables, mechanical stiffness boundaries, and 
				alignment of safety rules.
			\end{itemize}

	\subsection{Kinematic Gear Ratio Determination}
		The absolute target reduction ratio $u_{overall}$ is calculated from the input and output rotational velocities: 
			\begin{equation}
				\label{eq:reductionratio}
				u_{\text{overall}}
				\;=\;
				\frac{n_{\text{in}}}{\,n_{\text{out}}}
			\end{equation}
		This value defines the total kinematic reduction ratio. It establishes the exact multiplier by which the 
		rotational input speed drops, and the output torque increases. 

	\subsection{Design Torque Calculation}
		The opening stage of gearbox layout design demands calculating the nominal operating torque at the 
		output shaft ($T_{nominal}$) which is derived from the absorbed torque of the driven machine, or from the 
		nominal output torque of the prime mover ($T_{motor}$) scaled by the total reduction ratio ($u_{overall}$) 
		and mechanical efficiency ($\eta$):
			\begin{equation}
				\label{eq:nominaltorque}
				T_{\text{nominal}}
				\;=\;
				\frac{P}{\,n_{\text{in}}\!\Bigl(\frac{2\pi}{60}\Bigr)}
				\;\eta\;
				u_{\text{overall}}
			\end{equation}
		Where: 
			\begin{itemize}
				\item \textbf{$\eta$} represents the dimensionless mechanical efficiency coefficient of the gearbox. 
				This factor accounts for internal energy losses turned to heat, rolling resistance inside the shaft 
				bearings, and hydrodynamic churning resistance from the oil lubricant.
			\end{itemize}

		To safeguard against unpredictable real-world operating environments, this nominal torque value must be 
		multiplied by an empirical application (service) factor ($k_{a}$) and a life duty factor ($k_{t}$): 
			\begin{equation}
				\label{eq:designtorque}
				T_{\text{design}}
				\;=\;
				T_{\text{nominal}}\; k_a \; k_t
			\end{equation}

	\subsection{Preliminary Center Distance Estimation}
		Once the design torque is established, engineers estimate the initial center distance between gear axes 
		using empirical output torque curves: 
			\begin{equation}
				\label{eq:aestimate}
				\alpha_{\text{estimate}}  
				\;=\; 
				f(T_{\text{design}})
			\end{equation}
		This initial baseline geometry provides a starting point before running detailed strength calculations 
		under standards like AGMA or ISO.

	\subsection{Reduction Stage Determination}
		To preserve reliable gear tooth geometry and avoid overly slender pinions that twist or bend under load, the total 
		reduction ratio must be divided across an optimal number of intermediate reduction stages \autoref{tab:Table1}. 
			\begin{table}[ht]
				\renewcommand{\arraystretch}{1.5}
				\centering
				\caption{Ratio Range Selection}
				\label{tab:Table1}
				\small
				\begin{tabular}{|c|c|}
					\hline
        			\textbf{Number of Reductions} & \textbf{Economical Overall Ratio Range} \\					
					\hline
					Single Reduction   & Up to 6:1  \\
					\hline
					Double Reduction   & 5:1 to 25:1  \\
					\hline
					Triple Reduction   & 25:1 to 125:1  \\
					\hline
					Quadruple Reduction   & 125:1 to 625:1  \\
					\hline
				\end{tabular}
			\end{table}

		To distribute material costs and torque capacity efficiently, multi-stage gearboxes typically implement a larger reduction 
		ratio at the high-speed input stage, progressively decreasing the ratio toward the lower-speed output stages. 

		When splitting a total required reduction ratio ($u_{overall}$) across a multi-stage compound gear 
		train (such as a triple-stage reducer), handbook guidelines recommend establishing identical ratio sizing targets 
		across each stage. This balanced approach helps ensure uniform shafting diameters, uniform bearing sizing, and consistent 
		mounting methods throughout the assembly. 

		To determine the initial target ratio ($u_{stage}$) for each intermediate reduction step prior to tooth selection, 
		calculate the root of the overall ratio corresponding to the total number of reductions ($N_{stages}$):
			\begin{equation}
				\label{eq:initialtargetratio}
				u_{\text{stage}}
				\;=\;
				N_{\text{stages}}\;\sqrt{\,u_{\text{overall}}\,}
			\end{equation}

	\subsection{Pinion Tooth Selection and Undercutting Boundaries}
		The tooth count of the input pinion ($z_{1}$) cannot be chosen arbitrarily. To avoid undercutting during manufacturing—which 
		removes material from the tooth root base and severely weakens it—the engineer utilizes the analytical geometry formula 
		matching the selected helix angle ($\beta$) and normal pressure angle ($\alpha_{n}$). First, the transverse pressure angle 
		($\alpha_{t}$) is determined: 	
			\begin{equation}
				\label{eq:transversepressureangle}
				\alpha_{t}
				= \operatorname{\arctan}\!\left(
				\frac{\tan(\alpha_{n})}{\cos(\beta)}
				\right)
			\end{equation}

		Using the transverse pressure angle, the engineer solves for the absolute minimum number of teeth ($z_{\min}$) required to prevent undercutting: 
			\begin{equation}
				\label{eq:absoluteminnumofteeth}
				z_{\min}
				\;=\;
				\frac{2\,\cos(\beta)}{\sin^{\,2}\!\bigl(\alpha_t\bigr)}
			\end{equation}
		The engineer consults minimum tooth count reference charts and sets the pinion tooth selection to meet the safety margin. 

	\subsection{Parametric Module Trial and Standardization}
		Next, the framework defines the physical dimensions and structural pitch of the individual gear teeth, quantified as the 
		normal module ($m_{n}$). The engineer references minimum tooth count charts to estimate a pinion tooth number ($z_{1}$) that avoids undercutting 
		at specified speeds. They evaluate standard analytical formulas to yield a theoretical module: 
			\begin{equation}
				\label{eq:theoreticalmodule}
				m_{\text{theo}}
				\;=\;
				\left(\frac{2\,T_{\text{design}}}{\,z_{1}\;\psi\;\sigma_{b}\;Y}\right)^{1/3}
			\end{equation}
		Where ($\psi$) is the face width factor selected by the engineer from standard handbook look-up rows 
		and Y is the dimensionless Lewis form factor located manually via handbook tables matching the 
		chosen tooth count ($z_{1}$) and pressure angle ($\alpha_{n}$).
	
		Because manufacturing tooling is strictly constrained to specific integer intervals, the engineer must consult a 
		Preferred Metric Modules index. The calculated theoretical module is rounded upward to the nearest approved commercial designation.  
			\begin{equation}
				\label{eq:normalmodule}
				m_{\text{n}}
				\;=\;
				\operatorname{Standardized}\!\bigl(m_{\text{theo}}\bigr)
			\end{equation}
		Shifting to a standardized module alters the pitch diameters. The engineer must manually re-evaluate the entire 
		geometric stack to compute a modified center distance, completely overriding the initial chart-derived guess: 
			\begin{equation}
				\label{eq:angle}
				a \;=\;
				\frac{\,m_{\text{n}}\;z_{\text{1}}
				\;+\;m_{\text{n}}\;
				\bigl( z_{\text{1}}\;u_{\text{overall}} \bigr)}{2}
			\end{equation}
		Once the standardized module is locked, the true structural face width of the gear teeth is calculated: 
			\begin{equation}
				\label{eq:truestructuralfacewidth}
				b \;=\; \psi\,m_{\text{n}}
			\end{equation}

	\subsection{Spatial Boundary Audit and Interference Verification}
		The final phase of the pipeline overlays the calculated component 
		volumes (gear tip diameters, shaft dimensions, and raw bearing widths) directly onto the vehicle's structural 
		envelope. The engineer evaluates the absolute outermost tip diameter of the largest gear assemblies: 
			\begin{equation}
				\label{eq:absoluteoutermosttipdiameter}
				D_{\text{tip}}
				\;=\;
				m_{\text{n}}\;\bigl(z_{\text{2}} + 2\bigr)
			\end{equation}
		They physically verify if the rotating elements clear the static airframe frame and fit within 
		the designated housing box. 

	\subsection{AGMA Gear Tooth Bending Failure Analysis}
		Once the standardized geometry is finalized, the engineer evaluates the mechanical stresses at the tooth 
		contact points to ensure safety margins are met. To calculate the stress concentrations pulling at the 
		base of the tooth root, the tangential load ($F_{t}$) is determined: 
			\begin{equation}
				\label{eq:tangentialload}
				F_{\text{t}}
				\;=\;
				2 \cdot \frac{T_{\text{design}}}{d_{1}}
			\end{equation}
		The complete AGMA root bending stress ($\sigma_{b}$) equation incorporates dynamic adjustments for operational variables: 
			\begin{equation}
				\label{eq:agmarootbendingstress}
				\sigma_{b}
				\;=\;
				\frac{F_{\text{t}}\,K_{\text{a}}\,K_{\text{v}}\,K_{\text{m}}}
				{\,b\,m_{\text{n}}\,J}
			\end{equation}
		Where: 
			\begin{itemize}
				\item \textbf{$K_{a}$:} External application shock factor
				\item \textbf{$K_{v}$:} Dynamic velocity multiplier tracking tooth profiling errors
				\item \textbf{$K_{m}$:} Load distribution factor mapping shaft misalignment
				\item \textbf{b:} Structural face width of the gear mesh
				\item \textbf{J:} AGMA geometry factor for root bending
			\end{itemize}
		The safety criteria requires that the calculated stress remains below the material's allowable threshold. 

	\subsection{AGMA Surface Compressive Contact Fatigue Analysis}
		To protect the tooth profiles from compressive contact pitting, the Hertzian contact stress ($\sigma_{c}$)
		across the fluid film interface is calculated: 
			\begin{equation}
				\label{eq:hertziancontactstress}
				\sigma_{c}
				\;=\;
				Z_{\text{E}} 
				\,\sqrt{\frac{F_{\text{t}}\, K_{\text{a}}\, K_{\text{v}}\, K_{\text{m}}} {\,b\,d_{\text{1}}\,I}}
			\end{equation}
		Where: 
			\begin{itemize}
				\item \textbf{$Z_{E}$:} Elasticity coefficient matching the combined material moduli
				\item \textbf{I:} AGMA geometry factor for surface contact pitting resistance
			\end{itemize}

	\subsection{Dual-Shaft Fatigue Sizing via the ASME Elliptic Criterion}
		With the center distance locked, the engineer sizes the solid support shaft segments ($D_{1}$, $D_{2}$) 
		to resist catastrophic failure under combined loading conditions. The gear tooth loading vectors 
		generate a severe alternating bending moment (M) along the shaft, while the rotation transmits a 
		steady torsional shear torque (T). To calculate the absolute minimum safe shaft diameter ($d_{\min}$)
		required under a design safety factor (N), apply the ASME Elliptic Fatigue framework: 
			\begin{equation}
				\label{eq:absoluteminimumsafeshaftdiameter}
				d_{\min}
				\;=\;
				\left(\frac{32\,N}{\pi}\,\sqrt{\left(\frac{K_{\text{f}}\, M}{S_{\text{n}}}\right)^{\,2}
				\;+\;
				\frac{3}{4}\,\left(\frac{T}{S_{\text{y}}}\right)^{\,2}}
				\right)^{\!1/3}
			\end{equation}
		Where: 
			\begin{itemize}
				\item \textbf{$K_{f}$:} Fatigue stress concentration factor
			\end{itemize}

 		The manual shaft dimensioning must satisfy this sizing boundary: 
			\begin{equation}
				\label{eq:sizingboundary}
				D_{\text{shaft}} \;\ge\; d_{\min}
			\end{equation}

	\subsection{Sizing Limitations and Analytical Realities}
		These macro-sizing calculations provide only a preliminary estimation of the gearbox envelope under 
		idealized conditions. In practice, this represents a highly simplified starting point; an engineer 
		cannot rely solely on these broad formulas to finalize a design. Transitioning to a field-ready unit 
		requires extensive multi-physics iterations, as localized tooth stresses, dynamic misalignments, and 
		housing flex will fundamentally alter the system's real-world performance. A summary of the stages of 
		the gear design workflow is shown in \autoref{fig:gearbox2}.
			\begin{figure}[h]
				\centering
				\includegraphics[width=0.9\textwidth]{gearbox/gearbox2}
				\caption{Gear Design Workflow}
				\label{fig:gearbox2}
			\end{figure}
			\FloatBarrier

		While this baseline workflow suffices standard commercial gearboxes, it falls short in high-performance 
		sectors like aerospace, where weight and acoustic constraints are severe. To push past these classical 
		boundaries, specialized research initiatives are required to redefine standard optimization rules. A prime 
		example of this transition from standard commercial sizing to advanced drivetrain engineering is the 
		Advanced Rotorcraft Transmission (A.R.T.) program.